\chapter{Magnetic field: forces and sources}\label{ch:mag}

The electric field pushes along its lines; the magnetic one pushes
\emph{across} them --- $\mathbf F=q\mathbf v\times\mathbf B$. That cross
product is why magnetic fields curve paths without ever paying energy.

\section{Lorentz force: no work, all steering}\label{sec:mag-lor}

\begin{equation}
  \mathbf F=q\mathbf v\times\mathbf B.\label{eq:lor}
\end{equation}
Since $\mathbf F\perp\mathbf v$ identically, $P=\mathbf F\cdot\mathbf v=0$
--- \textbf{a static magnetic field does no work}, ever
(figure~\ref{fig:rh}). It only steers: every magnetic machine's real
work is paid elsewhere (electric field, mechanical source).

\begin{figure}[htbp]\centering
\begin{tikzpicture}[scale=0.9,font=\scriptsize]
\draw[->,very thick,accent] (0,0) -- (1.8,0) node[right]{$\mathbf v$};
\draw[->,very thick,amber] (0,0) -- (0,1.4) node[above]{$\mathbf F$};
\node at (0.9,-0.5) {$\mathbf B$ into page $\Rightarrow$ force up};
\foreach \x/\y in {0.7/1.0, 1.4/1.0}{\node at (\x,\y) {$\times$};}
\end{tikzpicture}

\caption{Right-hand rule: $\mathbf v\times\mathbf B$ for positive $q$.
For electrons, flip it.}
\label{fig:rh}
\end{figure}

\section{Circular and helical motion}\label{sec:mag-circle}

Uniform $\mathbf B$ plus perpendicular $\mathbf v$ gives a circle:
\begin{equation}
  r=\frac{mv}{qB},\qquad T=\frac{2\pi m}{qB},\label{eq:cyc}
\end{equation}
and the period is \emph{independent of speed} --- the cyclotron's whole
business model. Add a parallel component and the circle becomes a helix
(figure~\ref{fig:helix}).

\begin{figure}[htbp]\centering
\begin{tikzpicture}[scale=0.9,font=\scriptsize]
\draw[->,thick] (-2.5,0) -- (2.7,0) node[right]{$\mathbf B$};
\draw[very thick,accent,domain=-1.2:1.2,samples=80] plot (\x,{0.8*cos(360*\x)});
\draw[dashed] (-2.5,-0.8) -- (2.5,-0.8);
\draw[dashed] (-2.5,0.8) -- (2.5,0.8);
\node at (0,-1.5) {helix: circle $\perp\mathbf B$ + drift $\parallel\mathbf B$};
\node[align=left] at (0,-1.5) {};
\end{tikzpicture}

\caption{Helical motion: the perpendicular component circles, the
parallel one coasts.}
\label{fig:helix}
\end{figure}

\begin{example}[Cyclotron and mass spectrometer]\label{ex:cyc}
A proton ($m=1.67\times10^{-27}$\,kg) orbits in $B=1$\,T at
$10^{6}$\,m/s. Find $r$, $f$, and the speed a singly-charged ion of
twice the mass needs for the same $r$.
\par\smallskip
\textbf{Strategy.} Equation~\eqref{eq:cyc} directly.
\begin{solution}
$r=mv/qB=1.67\times10^{-27}\times10^{6}/(1.6\times10^{-19})=1.04$\,cm;
$f=qB/2\pi m=15.3$\,MHz --- RF range, why cyclotrons run RF.
Same $r$ with $2m$: $v=5\times10^{5}$\,m/s --- the mass spectrometer's
separation is just \eqref{eq:cyc} read as $m=rqB/v$. \checkmark
\end{solution}
\end{example}

\section{Force and torque on currents}\label{sec:mag-torque}

A wire segment feels $d\mathbf F=I\,d\boldsymbol\ell\times\mathbf B$;
a closed loop feels no net force in uniform $\mathbf B$ but a
\emph{torque} $\boldsymbol\tau=\boldsymbol\mu\times\mathbf B$ with
$\boldsymbol\mu=IA\hat{\mathbf n}$ (figure~\ref{fig:loop}) ---
the DC motor's entire principle in one vector product.

\begin{figure}[htbp]\centering
\begin{tikzpicture}[scale=0.95,font=\scriptsize]
\draw[very thick] (-1.3,-0.9) rectangle (1.3,0.9);
\draw[->,thick,amber] (-1.3,0) -- (-2.1,0) node[left]{$\mathbf F$};
\draw[->,thick,amber] (1.3,0) -- (2.1,0) node[right]{$\mathbf F$};
\foreach \x in {-0.6,0,0.6}{\node at (\x,0) {$\times$};}
\node at (0.9,-0.5) {$B$ in};
\draw[->,thick,accent] (0,1.3) -- (0.9,1.3) node[right]{$\boldsymbol\mu$};
\node[align=left] at (0,-1.6) {forces form a couple:\\$\tau=\mu B\sin\theta$};
\end{tikzpicture}

\caption{Loop in a field: forces couple into torque
$\boldsymbol\mu\times\mathbf B$, aligning $\boldsymbol\mu$ with
$\mathbf B$.}
\label{fig:loop}
\end{figure}

\begin{example}[Torque on a coil]\label{ex:torque}
A $100$-turn coil, $4\times6$\,cm, carries $2$\,A in $B=0.5$\,T at
$30^{\circ}$ to the field. Find the torque.
\par\smallskip
\textbf{Strategy.} $\tau=NIAB\sin\theta$.
\begin{solution}
$\tau=100\times2\times0.0024\times0.5\times\sin30^{\circ}=0.12$\,Nm ---
enough to spin a meter needle against a spring, which is precisely how
analog meters work. \checkmark
\end{solution}
\end{example}

\section{Hall effect: the carrier's signature}\label{sec:mag-hall}

Current through a field pushes carriers to one face
(figure~\ref{fig:hall}): the Hall voltage $V_H=IB/nqt$ builds until its
field cancels the magnetic push --- and its \emph{sign} reveals whether
the carriers are electrons or holes, the only cheap experiment that
answers that question.

\begin{figure}[htbp]\centering
\begin{tikzpicture}[scale=0.95,font=\scriptsize]
\draw[very thick] (-1.8,-0.7) rectangle (1.8,0.7);
\draw[->,very thick,accent] (-2.4,0) -- (-1.9,0) node[left]{$I$};
\foreach \x in {-1.2,-0.6,0,0.6,1.2}{\node at (\x,0) {$\times$};}
\node at (1.0,0.45) {$B$ in};
\draw[->,thick,amber] (0,-0.3) -- (0,0.35);
\node[align=left] at (0,-1.3) {carriers pushed to one face:\\$V_H$'s sign tells the carrier's sign};
\end{tikzpicture}

\caption{The Hall setup: transverse charge piles until $E_H=vB$; the
voltage's sign names the carrier.}
\label{fig:hall}
\end{figure}

\begin{example}[Hall voltage magnitude]\label{ex:hallv}
A copper strip ($n=8.5\times10^{28}$\,m$^{-3}$) of thickness
$t=1$\,mm carries $I=10$\,A across $B=1$\,T. Find $V_H$ and its sign
for electrons.
\par\smallskip
\textbf{Strategy.} $V_H=IB/nqt$.
\begin{solution}
$V_H=10\times1/(8.5\times10^{28}\times1.6\times10^{-19}\times10^{-3})
=7.4\times10^{-7}$\,V --- sub-microvolt: real Hall sensors use
semiconductors ($n\sim10^{22}$) for millivolt outputs. \checkmark
\end{solution}
\end{example}

\section{Biot--Savart: the source integral}\label{sec:mag-bs}

The field of a current element:
\begin{equation}
  d\mathbf B=\frac{\mu_0}{4\pi}\frac{I\,d\boldsymbol\ell\times\hat r}{r^{2}},
  \label{eq:bs}
\end{equation}
the magnetic counterpart of the Coulomb integral. For a long straight
wire it sums to $B=\mu_0I/2\pi r$ (figure~\ref{fig:wire}) --- same
$1/r$ as the line charge's field.

\begin{figure}[htbp]\centering
\begin{tikzpicture}[scale=0.9,font=\scriptsize]
\draw[very thick,accent] (0,-1.8) -- (0,1.8) node[above]{$I$};
\foreach \r in {0.5,1.0,1.5}{\draw[dashed] (0,0) circle (\r);}
\draw[->,thick,amber] (1.5,0) arc (0:70:1.5);
\draw[->,thick,amber] (-1.0,0) arc (180:110:1.0);
\node[align=left] at (0,-2.3) {$B=\mu_0I/2\pi r$ wraps the wire\\(right-hand grip rule)};
\end{tikzpicture}

\caption{Field of a straight wire: circles of strength
$\mu_0I/2\pi r$.}
\label{fig:wire}
\end{figure}

\begin{example}[Field on the axis of a current loop]\label{ex:loopaxis}
A circular loop of radius $R$ carries $I$. Find $\mathbf B$ on the
axis at distance $x$.
\par\smallskip
\textbf{Strategy.} Symmetry kills the transverse part; every element
contributes the same axial piece.
\begin{solution}
Each $d\mathbf B$ has axial factor $R/\sqrt{x^{2}+R^{2}}$:
\[B=\frac{\mu_0I}{4\pi}\frac{2\pi R^{2}}{(x^{2}+R^{2})^{3/2}}
=\frac{\mu_0IR^{2}}{2(x^{2}+R^{2})^{3/2}}.\]
Limits: centre gives $\mu_0I/2R$; far gives $\mu_0IA/2\pi x^{3}$
--- the magnetic dipole's $1/r^{3}$, matching the electric dipole of
chapter~9. \checkmark
\end{solution}
\end{example}

\begin{example}[The dipole's energy --- why magnets chase strong
fields]\label{ex:dipen}
A coil with $\boldsymbol\mu$ sits in $\mathbf B$ at angle $\theta$.
Find the energy needed to rotate it --- and the net force when
$\mathbf B$ is not uniform.
\par\smallskip
\textbf{Strategy.} Torque is $-dU/d\theta$: integrate
$\tau=\mu B\sin\theta$ against the angle. For the force, a dipole
that stays aligned does work on the field gradient.
\begin{solution}
$U(\theta)-U(0)=\int_0^{\theta}\mu B\sin\theta'\,d\theta'
=\mu B(1-\cos\theta)$, i.e.\ $U=-\boldsymbol\mu\cdot\mathbf B$:
stability lives at $\theta=0$, antiparallel is the hilltop. In a
non-uniform field the aligned dipole feels
$F_x=\mu\,\partial B/\partial x$ --- toward stronger $B$. A
$\mu=0.5$\,A\,m$^{2}$ coil in a gradient $0.1$\,T/m feels
$0.05$\,N, always in the direction of growing $|B|$. That is why
iron filings run to the poles, why magnetic separators work, and how
MRI's gradient coils drag spins into position: the dipole reads the
landscape of $|B|$ like a ball reads height. \checkmark
\end{solution}
\end{example}

\section{Amp\`ere's law: the magnetic Gauss}\label{sec:mag-amp}

\begin{equation}
  \oint\mathbf B\cdot d\boldsymbol\ell=\mu_0I_{\rm enc},\label{eq:amp}
\end{equation}
useful for the same three symmetries: wire ($2\pi r$), solenoid,
toroid. Limits: it fails off-symmetry exactly as Gauss does ---
chapter~10's warning reissued.

\begin{example}[Long solenoid by Amp\`ere]\label{ex:sol}
Find the field inside a solenoid of $n$ turns per length carrying $I$.
\par\smallskip
\textbf{Strategy.} Rectangular Amperian loop, one side inside
(figure~\ref{fig:sol}).
\begin{solution}
Outer sides give zero (field $\approx0$ outside); only the inner side
contributes $BL=\mu_0nLI$, so $B=\mu_0nI$ --- uniform, independent of
position inside. \checkmark
\end{solution}
\end{example}

\begin{figure}[htbp]\centering
\begin{tikzpicture}[scale=0.9,font=\scriptsize]
\foreach \x in {-1.6,-0.8,0,0.8,1.6}{\draw[very thick] (\x,0.6) ellipse (0.12 and 0.3);
\draw[very thick] (\x,-0.6) ellipse (0.12 and 0.3);}
\foreach \x in {-1.6,-0.8,0,0.8}{\draw[very thick] (\x,0.6) -- (\x+0.8,0.6);
\draw[very thick] (\x,-0.6) -- (\x+0.8,-0.6);}
\foreach \x in {-1.2,-0.4,0.4,1.2}{\draw[->,accent] (\x,0) -- (\x+0.35,0);}
\draw[->,accent] (-2.6,0) -- (-2.0,0);
\draw[->,accent] (2.2,0) .. controls (2.8,0.3) .. (2.6,0.9);
\node[align=left] at (0,-1.4) {inside: $B=\mu_0nI$ uniform;\\outside: $\approx0$};
\end{tikzpicture}

\caption{The solenoid's field: uniform inside, near-zero outside ---
Amp\`ere's gift.}
\label{fig:sol}
\end{figure}

\begin{lstlisting}[language=Python,caption={\texttt{code/ch13\_lorentz\_orbit.py} --- naive Euler vs Boris rotation in $q\mathbf v\times\mathbf B$.},label={lst:lor}]
vx, vy = c*vx + s*vy, -s*vx + c*vy   # exact rotation
x += vx*dt; y += vy*dt
\end{lstlisting}

Output of a real run:
\begin{lstlisting}[numbers=none]
theory: r=1.0, |v|=1.0 forever (B does no work)
naive Euler dt=0.1: |v|=19.788  r=19.788   <-- secular drift
naive Euler dt=0.01: |v|=1.350  r=1.350   <-- secular drift
Boris dt=0.1: |v|=1.0000  r in [0.9504,1.0504] (theory 1.0)
Boris dt=0.01: |v|=1.0000  r in [0.9950,1.0050] (theory 1.0)

helix pitch: vz=0.5 -> pitch=3.142 (period T=6.283 indep of v)
\end{lstlisting}

The honest result: plain Euler \emph{invents} energy ($|v|$ runs to
$19.8$) while the Boris rotation holds it at $1.0000$ and the radius
within a percent of theory --- the integrator must respect what the
force can't do.

\begin{example}[The Hall voltage: a voltmeter that weighs charge carriers]\label{ex:hall}
A copper strip $0.1$\,mm thick carries $10$\,A through a $0.5$\,T
field. Find the Hall voltage --- and why practical sensors are
semiconductor, not copper.
\par\smallskip
\textbf{Strategy.} Steady state: the transverse field balances
$\mathbf v\times\mathbf B$; with $I=nqv_dtA$, the Hall field is
$E_H=v_dB$, so $V_H=IB/(nqt)$.
\begin{solution}
$V_H=\dfrac{10\times0.5}{8.5\times10^{28}\times1.6\times10^{-19}
\times10^{-4}}=3.7$\,\textmu V for copper ($n=8.5\times10^{28}$/m$^{3}$).
A semiconductor at $n\sim10^{22}$/m$^{3}$ gives $10^{6}\times$ more
--- millivolts, a measurable signal. Copper's carriers are so dense
the Hall deflection is microscopic; dilute carriers deflect more.
\checkmark
\end{solution}
\end{example}

\begin{lstlisting}[language=Python,caption={\texttt{code/ch13\_two\_wire.py} --- wire--wire force, the old ampere, and a busbar's midline field.},label={lst:2wire}]
mu0 = 4 * np.pi * 1e-7
def force_per_m(I1, I2, d):
    return mu0 * I1 * I2 / (2 * np.pi * d)

print("   d (m)    F/L (N/m)   for I1=I2=100 A")
for d in [0.01, 0.05, 0.10, 0.50, 1.00]:
    print(f"{d:8.2f}   {force_per_m(100,100,d):10.4f}")
# the ampere definition check: solve force_per_m(I,I,1) = 2e-7
I = np.sqrt(2e-7 * 2 * np.pi * 1.0 / mu0)
print(f"I giving exactly 2e-7 N/m at d=1 m: {I:.4f} A")
# field midway between two antiparallel 100 A busbars 10 cm apart
d = 0.10
Bmid = 2 * mu0 * 100 / (2 * np.pi * (d / 2))
print(f"B midway between antiparallel 100 A bars 10 cm apart: {Bmid*1e6:.1f} uT")
\end{lstlisting}

Output of a real run:
\begin{lstlisting}[numbers=none]
   d (m)    F/L (N/m)   for I1=I2=100 A
    0.01       0.2000
    0.05       0.0400
    0.10       0.0200
    0.50       0.0040
    1.00       0.0020
I giving exactly 2e-7 N/m at d=1 m: 1.0000 A
B midway between antiparallel 100 A bars 10 cm apart: 800.0 uT
\end{lstlisting}

The force is small per metre --- a centimetre-separated 100\,A pair
pulls only $0.2$\,N/m --- until you remember fault currents: a
$20$\,kA short on the same spacing gives $8$\,kN/m, which is why
switchgear bracing is structural engineering. The middle line
recovering $I=1$\,A is not a coincidence: it \emph{was} the
definition (see workshop W25).

\section{In engineering practice: the magnetic toolkit}\label{sec:mag-eng}

Every DC motor is \S\ref{sec:mag-torque} with a commutator flipping
$\boldsymbol\mu$ at the right instant; mass spectrometers run
\eqref{eq:cyc} as a mass scale; Hall sensors price \S\ref{sec:mag-hall}
into current probes, throttle positions and phone compasses; and
MRI's solenoid is \S\ref{sec:mag-amp} stretched to 3\,T of uniform
field --- the chapter turned into hospitals.

\section{Common mistakes}\label{sec:mag-mistakes}

\begin{tripwire}[Left hand for a negative charge]
$\mathbf v\times\mathbf B$ is for positive $q$; electrons flip the
answer. Most wrong directions are one sign off.
\end{tripwire}
\begin{tripwire}[``$\mathbf B$ does work'']
It can't --- $\mathbf F\perp\mathbf v$ identically. Any energy change
in a magnetic problem was paid by an electric field or a battery.
\end{tripwire}
\begin{tripwire}[Amp\`ere off-symmetry]
$\oint\mathbf B\cdot d\boldsymbol\ell=\mu_0I$ always true, but
$B\cdot\oint d\ell$ factorisation needs the field constant on the loop
--- wire, solenoid, toroid only.
\end{tripwire}
\begin{tripwire}[$\mu_0$ vs permeability]
$\mu_0$ is vacuum; inside material use $\mu=\mu_r\mu_0$ and the
biot--savart/amp\`ere formulas pick it up.
\end{tripwire}
\begin{tripwire}[Helix pitch ignored]
The parallel speed doesn't circle --- the trajectory is a helix, and
``radius'' means the projected circle only.
\end{tripwire}

\section{Summary}\label{sec:mag-summary}

\begin{table}[ht]\centering
\begin{tabularx}{\linewidth}{@{}lY@{}}
\toprule
\tablehead\thead{Idea} & \thead{One-line version}\\
\midrule
Lorentz & $\mathbf F=q\mathbf v\times\mathbf B$; never works \\
orbit & $r=mv/qB$, $T=2\pi m/qB$ \\
torque & $\boldsymbol\tau=\boldsymbol\mu\times\mathbf B$, $\mu=IA\hat n$ \\
Hall & $V_H=IB/nqt$; sign = carrier sign \\
wire & $B=\mu_0I/2\pi r$ \\
solenoid & $B=\mu_0nI$ inside, $\approx0$ out \\
\bottomrule
\end{tabularx}
\caption{Magnetic field cheat sheet.}
\label{tab:magsummary}
\end{table}

\begin{wins}[Definition of done]
You can now: apply the right-hand rule both signs, compute cyclotron and
helical parameters, price loop torque, read a Hall voltage, integrate
Biot--Savart on an axis, and run Amp\`ere on wire/solenoid.
\end{wins}

\chapternotes{%
The force-first order follows MIT 8.02SC modules 16--20 \cite{ocw801}
and AP Physics C E\&M \cite{apced}; Purcell--Morin chs.~5--6 is the
canonical relativistic framing of where $\mathbf B$ comes from
\cite{hrw}. The integrator numbers are the listing's real output.}

\begin{rulebox}[{Watch instead}]
\begin{enumerate}
\item \emph{Magnetic Fields, Lorentz Force, Torques, Electric Motors
(DC)} --- Walter Lewin, MIT 8.02x Lecture 11, 49:50,
\url{https://youtu.be/0y9x7CS5Vrk} ---
\S\ref{sec:mag-lor}--\S\ref{sec:mag-torque} at full length (accessed
2026-10-11).
\item \emph{Force between two parallel current wires} --- Khan Academy,
11:07, \url{https://youtu.be/uxEB8xLrbCA} --- \S\ref{sec:mag-bs}
application worked quickly (accessed 2026-10-11).
\end{enumerate}
\end{rulebox}

\section*{Further reading}
\addcontentsline{toc}{section}{Further reading}

\begin{itemize}
\item MIT OCW \emph{8.02SC}, magnetic-field modules 16--20. At
\url{https://ocw.mit.edu}, course ``8.02SC'' (verified 2026-10-11).
\item OpenStax \emph{College Physics 2e}, ch.~22 ``Magnetism'' ---
free text: \url{https://openstax.org} (verified 2026-10-11).
\item College Board, \emph{AP Physics C (E\&M) CED}, Unit 12 ---
verified 2026-10-11 at \url{https://apstudents.collegeboard.org}.
\item E.~M. Purcell and D.~J. Morin, \emph{Electricity and Magnetism},
chs.~5--6 --- magnetism as relativity's corollary; the deeper anchor.
\item R.~P. Feynman, \emph{Lectures on Physics}, vol.~II, ch.~13 ---
the field view of magnetostatics; free at Caltech's site.
\item Khan Academy, \emph{Moving charges \& magnetism} --- drills
matching \S\ref{sec:mag-lor}--\S\ref{sec:mag-amp}, at
\url{https://www.khanacademy.org} (verified 2026-10-11).
\end{itemize}

\section*{Additional exercises}
\addcontentsline{toc}{section}{Additional exercises}

\begin{exercise}[easy]
An electron moves east in a field pointing up: which way does it
curve? (AP-C)
\end{exercise}
\begin{exercise}[easy]
A proton circles at $r=5$\,cm in $B=0.2$\,T: find its speed and
period. (AP-C)
\end{exercise}
\begin{exercise}[easy]
A $20$\,cm wire carries $5$\,A perpendicular to $B=0.3$\,T: find the
force. (AP-C)
\end{exercise}
\begin{exercise}[easy]
Find $B$ at $5$\,cm from a wire carrying $20$\,A. (F=ma)
\end{exercise}
\begin{exercise}[easy]
A loop's $\boldsymbol\mu$ is antiparallel to $\mathbf B$: describe the
torque and whether it's stable. (AP-C)
\end{exercise}
\begin{exercise}[medium]
Design a velocity selector: given $v=10^{6}$\,m/s and $B=0.01$\,T,
find $E$ for undeflected passage. (AP-C FR-style)
\end{exercise}
\begin{exercise}[medium]
A $50$-turn $10$\,cm square coil in $0.8$\,T carries $3$\,A: find
maximum torque and where it occurs. (AP-C)
\end{exercise}
\begin{exercise}[medium]
Two parallel wires $10$\,cm apart carry $20$\,A the same way: find
force per metre and direction; flip one and redo. (F=ma)
\end{exercise}
\begin{exercise}[medium]
Explain how a mass spectrometer separates isotopes given fixed $B$,
$v$: what quantities differ? (F=ma)
\end{exercise}
\begin{exercise}[medium]
A coaxial cable's Amp\`erian field: $B(r)$ inside inner conductor
(uniform current), between, and outside. (engineering)
\end{exercise}
\begin{exercise}[medium]
Estimate the Hall voltage for a $1$\,mm-thick copper strip at $5$\,A
in $0.5$\,T. (F=ma)
\end{exercise}
\begin{exercise}[hard]
In crossed $\mathbf E$ and $\mathbf B$ ($E=v_dB$ drift): describe the
cycloid trajectory and its drift velocity. (USAPhO-style)
\end{exercise}
\begin{exercise}[hard]
A finite wire of length $2L$: find $B$ at perpendicular distance $r$
and recover the infinite-wire limit. (APhO-style)
\end{exercise}
\begin{exercise}[medium]
A toroid of $N$ turns, radius $R$: find $B$ inside the ring by
Amp\`ere and explain why it depends on $r$. (AP-C)
\end{exercise}
\begin{exercise}[medium]
An electron enters a field at $45^{\circ}$ to $\mathbf B$: describe
the helix's radius and pitch. (F=ma)
\end{exercise}
\begin{exercise}[easy]
A proton travels parallel to a current-carrying wire, moving the same
way as the conventional current: find the force's direction on it.
(AP-C)
\end{exercise}
\begin{exercise}[medium]
A length of wire is bent into one circular turn of radius $R$, then
into two stacked turns of radius $R/2$: compare $B$ at the centre.
(AP-C)
\end{exercise}
\begin{exercise}[medium]
A long solenoid's field is $\mu_0nI$ deep inside: argue why the field
at the open end is about half that, and where the ``other half''
lives. (F=ma)
\end{exercise}
\begin{exercise}[hard]
A dipole $\boldsymbol\mu$ is \emph{held} antiparallel to a field
$B(x)$ that grows with $x$: find the force on it, and explain why
releasing it flips the answer. (APhO-style)
\end{exercise}
